% ===========================================================================
%  hAIthem blog post, editable source
%
%  Write the prose here; hand the file back and it gets pushed to
%  blog/bsm-rl-search.html. Everything the page renders has a placeholder
%  below, in page order, so moving a \interactive or \figplace block moves
%  the real thing on the page.
%
%  Compile with pdflatex to see the layout with the placeholders drawn in.
%
%  Rules of the mapping:
%    \section{...}        -> an <h2> on the page
%    \subsection{...}     -> an <h3> on the page
%    ordinary paragraphs  -> <p>
%    \lead{...}           -> the opening paragraph, set larger
%    \interactive{key}{}  -> a live figure, keys listed below
%    \figplace{file}{}{}  -> a static figure from haithem-assets/figures
%    figrow environment   -> the figures inside share one row on the page
%    \pullquote{...}      -> a large centred quote between sections
%    \pullquote*{...}     -> the same panel without the quotation marks
%    \pullquote[b]{a}     -> two of them, "a" then an arrow then "b"
%    \wide                -> as the last argument, breaks the figure out of
%                            the text column (for very wide diagrams)
%    \haithem \alphago    -> names, styled the same wherever they appear
%    \alphazero
%    thebibliography      -> the References list at the foot of the page,
%                            \cite keys become superscript links into it
%
%  Interactive keys, do not rename:
%    alphazero  the bundled AlphaZero parameter-space page
%    param3d    the 3D parameter-space explorer
%    tree       the constraint-tree explorer (walk + whole tree)
% ===========================================================================

\documentclass[11pt]{article}
\usepackage[margin=1in]{geometry}
\usepackage{xcolor}
\usepackage{tcolorbox}
\usepackage{hyperref}
\usepackage{parskip}

% names that get the same bold small-caps treatment wherever they appear
\newcommand{\haithem}{\textbf{\textsc{hAIthem}}}
\newcommand{\alphago}{\textbf{\textsc{AlphaGo}}}
\newcommand{\alphazero}{\textbf{\textsc{AlphaZero}}}

\newcommand{\posttitle}[1]{\title{#1}}
\newcommand{\posttags}[1]{\def\theposttags{#1}}
\newcommand{\postdate}[1]{\date{#1}}
\newcommand{\postbyline}[1]{\def\thepostbyline{#1}}

% opening paragraph, rendered larger on the page
\newcommand{\lead}[1]{{\large #1}\par\medskip}

% a live figure on the page
\newcommand{\interactive}[3][]{%
  \begin{tcolorbox}[colback=orange!5,colframe=orange!60!black,
                    title={Interactive: \texttt{#2}\ifx\relax#1\relax\else\ (#1)\fi}]
  \emph{Caption:} #3
  \end{tcolorbox}}

% a static figure; #1 optional "wide", #2 file stem, #3 caption
\newcommand{\figplace}[3][]{%
  \begin{tcolorbox}[colback=blue!4,colframe=blue!40!black,
                    title={Figure: \texttt{#2}\ifx\relax#1\relax\else\ (#1)\fi}]
  \emph{Caption:} #3
  \end{tcolorbox}}

% the figures inside share one row on the page; on paper they simply follow
% one another
\newenvironment{figrow}{}{}

% a large centred quote. \pullquote{one} on its own, or \pullquote[after]{before}
% for the two-part form with an arrow between them
\makeatletter
\newcommand{\pullquote}{\@ifstar\pq@bare\pq@quoted}
\newcommand{\pq@quoted}[2][]{%
  \begin{center}\itshape\large
  ``#2''\ifx\relax#1\relax\else\ $\Rightarrow$\ ``#1''\fi
  \end{center}}
\newcommand{\pq@bare}[2][]{%
  \begin{center}\itshape\large
  #2\ifx\relax#1\relax\else\ $\Rightarrow$\ #1\fi
  \end{center}}
\makeatother

% the downloads picker at the foot of the post
\newcommand{\datalist}[1]{%
  \begin{tcolorbox}[colback=gray!6,colframe=gray!50,
                    title={Generated: downloads picker}]
  #1
  \end{tcolorbox}}

\posttitle{The Search for (Non-Excluded) Beyond Standard Model Physics is an RL Problem}
\postdate{September 2026}
\posttags{Large Lagrangian Models; Dark Matter; Reinforcement Learning; LLM Agents}
\postbyline{by Ibrahim Elsharkawy, last edited September 2026}

\begin{document}
\maketitle
\noindent\textbf{Tags:} \theposttags

% ---------------------------------------------------------------- opening --
\interactive{alphazero}{Illustration of the similarities between \alphago{}
and \alphazero{} and \haithem{}. Both search optimal moves conditioned on
the current play turn and past actions.}

% the byline sits under the opening figure, not in the page header
\noindent\textit{\thepostbyline}\par\medskip


The current state of the search for beyond standard model (BSM) physics is one of experimental drought and theoretical abundance. No new elementary particle has been confirmed since the Higgs in 2012, and many dedicated experiments across high-energy physics (including the Large Hadron Collider) have found little significant deviation from the Standard Model across thousands of observables. Despite this, over the same period, the number of BSM descriptions has grown significantly \cite{Bertone:2018krk}. 

This is to say, there is an overabundance of ways to describe reality, but not enough experimental probes that can tell these scenarios apart. This state of physics thus motivates asking (and potentially automating a search for the answer of) a different question,  

% \pullquote*[What \emph{set of descriptions} are consistent with current data and
% \emph{what unexplored observables} distinguish them?]{What is the best description of
% reality?}

\pullquote*[What are unexplored \textbf{observables} that distinguish these descriptions?]{What is the best \textbf{description} of
reality?}


We introduce \haithem{} (\textbf{H}idden-sector \textbf{AI} for \textbf{T}estable \textbf{H}ypothesis \textbf{E}nu\textbf{M}eration), a framework designed to search the space of experimental signatures beyond the literature, backronymed after a pioneer of the experimental method, Ibn Al-Haytham (965-1040 CE), whose work cemented controlled experiment as the arbiter of competing theory \cite{Sabra:1989optics,Ibn-al-haythem-reflect}. \haithem{} turns the task of signature generation into a competition between candidate signatures' ability and practicality for separating two or more descriptions of reality that satisfy current experimental constraints.

\interactive{param3d}{Watch a trained Large Lagrangian Model play the
battleship-like search over parameter space.}

To do so, \haithem{} is first composed of what we call a Large Lagrangian Model (LLaM), an autoregressive transformer pretrained from scratch on over 1 Billion Output Tokens from over $10,000$ specially tokenized Lagrangians, and fine-tuned in a live environment with reinforcement learning (RL). The LLaM is trained to play a battleship-style game (with limited turns and limited probes per turn) over the parameter space of a theory. Its goal is to find regions in that parameter space that can be the correct description of reality (that is, no experiment we consider rules them out).

For a given Lagrangian, this is similar in structure to \textsc{AlphaGo} and \textsc{AlphaZero} \cite{silvergo,SilverGo2,doi:10.1126/science.aar6404} and our approach borrows concepts from each of them. We have a space too large for systematic search, made tractable by an RL agent trained with a physically grounded “game” engine. We find the most effective method to build such an agent is to adopt the paradigm that defines modern large language models and \textsc{AlphaGo}. That is, large-scale pretraining on an offline dataset, followed by Reinforcement Learning fine-tuning in a live environment.

\interactive{tree}{\textbf{Walk the tree:} Answer what observables were found
and find what Lagrangians survive. \textbf{Whole tree:} show the full tree,
scroll to zoom, and click a node for its details.}

With a set of non-excluded regions found by the LLaM, \haithem{} constructs a decision tree that generates the competition between candidate signatures. The decision tree formulation coupled to the LLaM is built to encourage the search for signatures beyond the literature. Each node in the tree is an observable with a binary outcome that splits LLaM-found regions (for example, do we expect an observation at experiment X or not), and the leaves of the tree correspond to parameter-space regions that share an experimental signature given by the set of nodes that comprise the path to the leaf. We first construct the tree with analytically computed observables, depth-ordering observables by their ability to split the set of regions we consider. In this construction, degenerate leaves are leaves with more than one parameter-space region sharing every experimental signature we compute analytically. These leaves are then handed to a set of LLM agents that attempt to propose realistic observables that continue the splitting. 

The tree construction turns breaking these degeneracies into a concrete, non-open-ended task given specific Lagrangians and specific parameter values (masses, couplings, etc.). The LLM-agents must compete to give realistic signatures along with sharp binary verdicts as to how degenerate regions differ in their proposal. Additionally, they are forced to think beyond observables already analytically computed earlier in the tree (reuse is explicitly prohibited) and are allowed null answers (no signature found). We further find that the decision tree formulation leads to a natural grouping of Lagrangian-viable regions into experimentally distinguishable classes.

% ----------------------------------------------------------------- method --
\section{The Large Lagrangian Model}


% \figplace{fig-lagrangian-tokens}{Tokenization of an example Lagrangian. Each
% dark-sector field is embedded into its own field token, while the symmetry
% content and per-field summaries (dashed) are embedded into a single global
% Lagrangian token.}

\figplace[wide]{fig-token-architecture-pretrain}{The Large Lagrangian Model.}
We let the search for regions in high-dimensional parameter space not excluded by experiment be akin to a single-player game of battleship. The board is defined by a Lagrangian, some set of experimental
cuts, the parameter space ranges, and a budget of turns the agent has to do the search, with a fixed number of “shots” per turn. The ships the LLaM looks for are connected regions that are both \textit{viable}, in that they are not ruled out, and \textit{testable}, in that some experiment is predicted to probe them. Their number, shape, and location are unknown, and the agent is rewarded in training for finding them efficiently.

The Large Lagrangian Model is trained to learn how to play this game and is an autoregressive (in parameter space dimension) transformer whose tokens are Lagrangians and their phenomenology.  Each episode, it is given one Lagrangian and a budget of between 5 and 50 turns. 
On every turn it proposes four distributions over that theory's parameter space, and the shots for that turn are drawn from them. The outcome of the current turn is given to the LLaM for it to use for the next turn(s).

The model takes a few inputs: the tokenized Lagrangian, the viability configuration (which experimental cuts are active), what ranges the agent can search in, and how many turns it has left. It also takes an encoding of the search so far, where 64 learned tokens
cross-attend over every probe shot in the game. This lets each proposal depend on what the previous turns revealed.

Each of the four proposed distributions is a Beta distribution per parameter. Independent Betas cannot express correlations between parameters, so
we sample them autoregressively, with masses first, then couplings, then kinetic mixing, with each drawn value fed back causally for the next parameter. 

Training in a live environment alone proved intractable as the model grew. We therefore pretrained offline on roughly a billion output tokens drawn from 10,000 distinct Lagrangians. The pretraining task is a perturbation of the real RL game, and we fine-tune with Proximal Policy Optimization (PPO) in a live environment. We build and train two sizes of the LLaM, one 4.4M (small) and one 42.5M parameters (medium).
% \figplace[wide]{fig-ar-sampling}{A depiction of autoregressive sampling. Each
% parameter is drawn from a model forward pass and fed back as context for the
% next forward pass for the next parameter. The mode chooses the beta
% distribution or differential evolution style policy once per probe and fixes
% the choice for all parameter dimensions.}

% ---------------------------------------------------------------- results --
\subsection{Benchmarking the Large Lagrangian Model}
Below are benchmarks of the LLaM against a Differential Evolution (DE) baseline. The LLaM finds more viable points than DE, especially with a small number of probes. The LLaM also does well with increasing Lagrangian parameter dimension, while DE collapses, as is expected with the curse of dimensionality.


\begin{figrow}
\figplace{total-viable-count-vs-budget-summed-over-21-benchmark-lagrangians}{%
Total number of viable points found as a function of search budget, summed
over 21 benchmark Lagrangians.}

\figplace{mean-viable-per-lagrangian}{Mean number of viable points per
benchmark Lagrangian as a function of the Lagrangian parameter-space dimension
$d$. The LLaM becomes \textit{more} effective as the dimension grows, while
Differential Evolution collapses, as is expected with the curse of
dimensionality.}
\end{figrow}

% \figplace[wide]{diversity-benchmark}{Diversity benchmark. Distinct
% experimental-signature classes discovered (that is, the number of decision
% tree leaf nodes, solid, left axis) and viable regions found per 100 viable
% points.}

\section{Building the Decision Tree}

\pullquote*{The Large Lagrangian Model finds \emph{what theories to consider}, and the \emph{LLM-agent does the phenomenology there}, both working at scale.}

We take all viable points the LLaM found for a Lagrangian, and
grow a tree with them, where each node is one observable, and each leaf defines a distinct combination of signatures and contains all viable points that share the same signatures. After growing the tree, we find leaves holding more than one region. We consider these degenerate, in that nothing we currently compute tells the theories in the leaf apart. 

Breaking these degeneracies is the job of LLM-agents, which can reason at scale.  We hand these leaves to an LLM agent and ask it to attempt to split the degeneracy, one LLM-agent call per leaf. The agent is asked to extend the tree with new binary splits until either all degeneracies have been split, a max LLM-depth is reached (we use a max LLM-depth of 4), or a proposed node is judged to be practically impossible to split. For adding binary splits, the agent has three choices for types of splits:
\begin{enumerate}
    \item A \textbf{Literature Split}: A published result for an observable whose current limits split the degenerate regions. For this kind of split, the LLM must cite the paper that is used, its reasoning for why this works, and a status \textsc{``Splits!"} or \textsc{``No Split"} that reports if the split was successful.
    \item A \textbf{Literature Projection}: An experiment or limit that is planned or described in the literature whose projected sensitivity splits the regions. For this split, the LLM must cite the most relevant papers for this projection, must report its reasoning for this split, an approximate sensitivity improvement factor relative to the closest experiment already run, and must report a qualitative judgment on the feasibility of this split as a choice from five categories: \textsc{Reanalysis}, \textsc{Possible}, \textsc{Next Generation}, \textsc{Speculative}, and \textsc{Impossible}.
    \item A \textbf{Novel Observable}: An experiment or limit that has not been proposed in the literature that may be capable of splitting the regions. The LLM agent must report its reasoning for this split, the paper that is closest to its idea, and must report a qualitative judgment on the feasibility of this split as a choice from four categories: \textsc{Possible}, \textsc{New Experiment}, \textsc{Speculative}, and \textsc{Impossible}.
\end{enumerate}

We enforce ordering rules for the proposal of these node choices. The first attempted split must be of type Literature Split to ensure the LLM searches what has already been done, and only after finding \textsc{``No Split"} is it permitted to propose other types. The agent then searches for \textbf{Literature Projection} nodes, and moves on to \textbf{Novel Observable} only if none are found. We also allow the LLM to propose alternatives. If a \textbf{Literature Projection} is qualitatively judged to be \textsc{Next Generation} or worse, the agent can propose a \textbf{Novel Observable} alternative with the same or better feasibility qualitative judgment.

We further use ensembling, creating the full tree $N_{\rm repeat} = 5$ times and passing all trees through an aggregation LLM-agent pass, in which the agent chooses the best combination of observables from all runs. For all runs, tool use is enabled and set to web search and arXiv retrieval, which the agent must use to verify every reference it cites. 


\section{What did we find?}
\figplace[wide]{paleo-detector}{The agent's proposal, a gypsum paleo-detector. The
sample is buried for a billion years and particles interact with the crystal lattice. The ratio of the longest calcium and oxygen tracks tells us the
dark matter mass independent of the galactic escape velocity.}
As an initial test, we use our LLaM to scan through all single-dark-multiplet models our space allows and build our decision tree over the (small!) set. We find that degenerate leaves lead to the proposal of unapplied combinations of known observables.

A specific example is a ``Novel Alternative" titled \textit{``Paleo-detector Ca/O recoil-spectrum ratio"}. The agent is given two models indistinguishable by all experiments in tree and differ only by their dark matter mass. The agent recognizes that dark matter mass can be probed by looking at maximum energy deposition seen in direct detection experiments, which is a function of the mass and the maximum expected velocity of the particle. This maximum velocity is given by the galactic escape velocity. 

The agent identifies two issues with the approach. First, this requires many detection events seen near the maximum velocity, which in turn requires a very large time (exposure) of the direct detection experiments, and secondly, the galactic escape velocity is a source of error. As a solution to the first, it proposes the use of paleo detectors. These are minerals found deep underground, protected from cosmic rays and high temperatures for billions of years. Disturbances of the crystal lattice of these detectors contain permanent records of nuclear recoil events throughout their lifetime from many sources of background, along with potential dark matter interactions. For the second issue, the agent proposes the use of a ratio of the maximum energy deposition for two different nuclei, and specifically proposes the calcium and oxygen in gypsum. The ratio of the two remains a function of the dark-matter mass, but the dependence on the galactic escape velocity cancels. 

Both paleo detectors and the ratio of observables have been proposed before, but the combination for this context appears to be unexplored. The proposal requires an in-depth analysis of its realism, but its existence motivates \haithem{}.


\section{Next Steps}

This work opens the door to large-scale searches for unexplored observables across different classes of dark matter and broader BSM Lagrangians. Since combinatorics grows quickly, such scans would require a large quantity of compute and potentially intelligent Lagrangian sampling strategies. The payoff, however, may be substantial. Such searches may reveal the value of different future and proposed observables/experiments in distinguishing between diverse dark matter and BSM models, and may help the community make better-informed scientific investments. 

% ----------------------------------------------------------- on the shelf --
% Additional materials is commented out on the page as well. Uncomment here
% and there together to bring the downloads picker back.
%
% \section{Additional materials}
%
% Pick a Lagrangian class for its decision tree at full resolution, and for
% the raw replies the literature-search agents gave when they proposed the
% splits in it.
%
% \datalist{Built at page load from
% \texttt{haithem-assets/downloads/index.json}: a dropdown grouped by agent
% run, and two buttons under it, the tree as a PDF and the agent transcripts
% as one Markdown file. Each button greys out where that run has nothing to
% give. Nothing to write here, it updates itself when the downloads are
% rebuilt.}

% ------------------------------------------------------------ references --
\begin{thebibliography}{9}

\bibitem{Bertone:2018krk}
G.~Bertone and T.~M.~P.~Tait, \emph{A new era in the search for dark matter},
Nature \textbf{562} (2018) 51--56,
\href{https://arxiv.org/abs/1810.01668}{arXiv:1810.01668}.

\bibitem{Sabra:1989optics}
Ibn al-Haytham, \emph{The Optics of Ibn al-Haytham, Books I--III: On Direct
Vision}, translated with introduction and commentary by A.~I.~Sabra, 2 vols.,
Studies of the Warburg Institute Vol.~40, Warburg Institute, University of
London, London (1989). ISBN 0-85481-072-2.

\bibitem{Ibn-al-haythem-reflect}
J.~Al-Khalili, \emph{In retrospect: Book of Optics}, Nature \textbf{518}
(2015) 164--165, \href{https://doi.org/10.1038/518164a}{doi:10.1038/518164a}.

\bibitem{silvergo}
D.~Silver \emph{et al.}, \emph{Mastering the game of Go with deep neural
networks and tree search}, Nature \textbf{529} (2016) 484--489,
\href{https://doi.org/10.1038/nature16961}{doi:10.1038/nature16961}.

\bibitem{SilverGo2}
D.~Silver \emph{et al.}, \emph{Mastering the game of Go without human
knowledge}, Nature \textbf{550} (2017) 354--359,
\href{https://doi.org/10.1038/nature24270}{doi:10.1038/nature24270}.

\bibitem{doi:10.1126/science.aar6404}
D.~Silver \emph{et al.}, \emph{A general reinforcement learning algorithm that
masters chess, shogi, and Go through self-play}, Science \textbf{362} (2018)
1140--1144,
\href{https://doi.org/10.1126/science.aar6404}{doi:10.1126/science.aar6404}.

\end{thebibliography}

\end{document}
